\pagenumbering{roman}%
\begin{abstract}%
    Los lahares constituyen uno de los principales peligros asociados a volcanes cubiertos por nieve y hielo, debido a su capacidad para movilizar rápidamente mezclas de agua, sedimentos y material volcánico a través de los 
    sistemas de drenaje. El volcán Villarrica, uno de los sistemas volcánicos más activos de Chile, constituye un caso de estudio relevante debido a la generación de lahares durante su evento eruptivo de 2015, asociados a la 
    interacción del proceso eruptivo con la cubierta de nieve y hielo.

    Esta investigación propone desarrollar un modelo computacional de lahares mediante Redes Neuronales Informadas por la Física (PINNs), utilizando el evento eruptivo de 2015 del volcán Villarrica como caso de estudio. 
    El modelo físico se formulará mediante un sistema acoplado de ecuaciones diferenciales parciales basado en la aproximación de flujo somero, incorporando la conservación de masa y momento, el transporte de concentración 
    sólida y la evolución de la presión de poros. La PINN aproximará los campos espacio-temporales de espesor, velocidad, concentración sólida y presión de poros, incorporando las ecuaciones gobernantes directamente en su 
    función de pérdida.

    La metodología considerará condiciones iniciales y de frontera físicamente consistentes, representación de la topografía y estrategias específicas para mejorar la estabilidad del entrenamiento, incluyendo la imposición 
    fuerte de condiciones iniciales, normalización de variables y control de la convergencia hacia soluciones triviales. Posteriormente, el modelo será aplicado al caso del Villarrica 2015 utilizando parámetros y observaciones 
    disponibles para el evento. Se espera evaluar la capacidad de las PINNs para reproducir la dinámica espacio-temporal de lahares y establecer su potencial como herramienta complementaria para la modelación 
    física de este tipo de peligro volcánico.
    
    %Los lahares representan uno de los peligros volcánicos más devastadores debido a su alta movilidad y potencial de impacto en comunidades e infraestructura. En Chile, 
    %el volcán Villarrica constituye un escenario crítico para el estudio de estos flujos, particularmente aquellos originados por la interacción entre magma y criósfera, 
    %como se observó en el evento eruptivo de 2015. Este trabajo presenta el desarrollo de un modelo computacional para la simulación de lahares basado en 
    %Redes Neuronales Informadas por la Física (PINNs). La arquitectura propuesta integra ecuaciones diferenciales parciales acopladas bajo la aproximación de flujo somero, 
    %incorporando variables de transporte de concentración sólida y dinámica de presión de poros para capturar la naturaleza multifásica del fenómeno   
    %La metodología aborda la construcción de una PINN entrenada mediante restricciones físicas y condiciones iniciales derivadas de datos históricos del evento de 2015. 
    %Un componente central del estudio consiste en el análisis y mitigación de patologías propias del entrenamiento de PINNs en sistemas geofísicos complejos, tales como 
    %la convergencia a soluciones triviales, el desbalance de gradientes y la inestabilidad numérica. Los resultados buscan demostrar la viabilidad de las PINNs como una
    %herramienta robusta para la modelación de peligros naturales, integrando las leyes de la física con los avances en inteligencia artificial para fortalecer las capacidades de predicción 
    %y gestión de riesgos volcánicos. 

\end{abstract}

\textbf{Palabras clave: Redes Neuronales Informadas por la Física (PINNs), Ecuaciones Diferenciales Parciales, Volcán Villarica, Lahares, Peligros Volcánicos} \palabrasClave\par%
